\subsection{The politico-economic equilibrium with CRRA}\label{\refsec:PEE}
With CRRA preferences, the past savings level $s_{t-1}$ becomes a relevant state variable for the political problem, and with it the decoupling of \eqref{\refmodeleq:PEELOG:decoupling} is lost: the continuation policy re-enters through $\bm{B}_{t+1}$ and $\Gamma_{s,t}$, and we are back to a backward recursion over policy \emph{functions} $\tau_t(s_{t-1})$. The state-space reduction survives otherwise intact (the savings \emph{shares} $s_{t-1,i}/s_{t-1}$ remain closed form in $\tau_t$, so the single aggregate level is the only state), and the machinery of \cref{\refsec:RobustRoot} applies along the policy dimension unchanged. What is new is the grid structure over the state, and a fixed-point problem that pins the endogenous state down.

\smalltitle{Grids.} Let $\mathcal{S} = \left\lbrace s^{(1)}, ..., s^{(M_s)} \right\rbrace$ with $s^{(1)} = l_s>0$ and $s^{(M_s)} = u_s>l_s$ denote a grid of states. Realistic ranges come from steady state dynamics: $s^{\star}(\tau)$ is monotonically decreasing in $\tau$ with $s^{\star} \rightarrow 0$ as $\tau \rightarrow 1$, so $l_s$ is a low value serving numerical stability and the upper bound is taken as $u_s = 1.25 \times s^{\star}(0)$. $\mathcal{S}$ is a grid in \emph{levels}, so by \eqref{\refmodeleq:scaleInvariance} its placement, unlike the solution it supports, depends on the units convention of the calibration, and by default it is rebuilt at every solve from the steady state at the current parameters. Let $\mathcal{S}\times\mathcal{G}_{\tau}$ denote the Cartesian grid of states and taxes.

For $t<T$ we additionally need a grid $\mathcal{S}^{\prime}$ of \emph{candidate} values for the endogenous state $s_t$: $\mathcal{S}$ indexes the predetermined states we solve \emph{at}, while $\mathcal{S}^{\prime}$ is searched \emph{over} to locate an equilibrium, and although the two may hold the same values, the accuracy of the search below is governed by $\mathcal{S}^{\prime}$ alone.

\smalltitle{The economic equilibrium at $t<T$ is a fixed point in $s_t$.}
Suppose we have solved period $t+1$ and hold continuation policies $\tau^{t+1}(\cdot)$ and $h^{t+1}(\cdot)$, both functions of the state $s_t$ entering $t+1$. For a candidate current tax $\tau_t$ and a candidate state $s_t$, the auxiliary equations \eqref{\refeq:auxiliary} can be evaluated in one forward pass:
\begin{subequations}\label{\refeq:stateApprox}
\begin{align}
    \bm{B}_{t+1} &= \bm{B}_{t+1}\left(s_t,\, h^{t+1}(s_t)\right), \label{\refeq:stateApprox:B}\\
    \Gamma_{s,t} &= \Gamma_{s,t}\left(\bm{B}_{t+1},\, \tau^{t+1}(s_t),\, \theta_{t+1}\right), \label{\refeq:stateApprox:Gammas}\\
    \Theta_{h,t} &= \Theta_{h,t}\left(\tau_t,\, \tau^{t+1}(s_t),\, \theta_{t+1},\, \Gamma_{s,t}\right), \label{\refeq:stateApprox:Thetah}\\
    \Theta_{s,t} &= \Theta_{s,t}\left(\Theta_{h,t},\, \Gamma_{s,t}\right). \label{\refeq:stateApprox:Thetas}
\end{align}
\end{subequations}
The candidate $s_t$ is an equilibrium at $s_{t-1}$ precisely when it reproduces itself through \eqref{\refeq:auxiliary:s}. Defining the residual
\begin{align}\label{\refeq:stateResidual}
    \mathcal{R}_t\left(\tau_t, s_t;\, s_{t-1}\right) \equiv \Theta_{s,t}\left(\tau_t, s_t\right)\left(\frac{s_{t-1}}{\nu_t}\right)^{\frac{\alpha(1+\xi)}{1+\alpha\xi}}-s_t,
\end{align}
the state policy function $s_t(\tau_t, s_{t-1})$ is the root of \eqref{\refeq:stateResidual} in its second argument. This is a genuine root-finding problem in $s_t$, not a recursion we can unroll, because $s_t$ enters the right hand side through $\bm{B}_{t+1}$ and the continuation policies. It is a \emph{root} problem, not a maximisation: any crossing is a solution, so neither the equilibrium test nor the selection of \cref{\refsec:RobustRoot} applies. Where the residual has several crossings on $\mathcal{S}^{\prime}$, the lowest is taken and the number of roots is reported.

\smalltitle{The fixed point is two-dimensional, not three.}
\emph{None} of $\bm{B}_{t+1}$, $\Gamma_{s,t}$, $\Theta_{h,t}$, $\Theta_{s,t}$ in \eqref{\refeq:stateApprox} depends on $s_{t-1}$: the predetermined state enters \eqref{\refeq:stateResidual} only through the explicit factor $(s_{t-1}/\nu_t)^{\alpha(1+\xi)/(1+\alpha\xi)}$. The costly part of the state approximation is therefore a function of $(\tau_t, s_t)$ alone, evaluated once on $\mathcal{G}_{\tau}\times\mathcal{S}^{\prime}$, after which the residual for \emph{every} $s_{t-1}\in\mathcal{S}$ follows by multiplying $\Theta_{s,t}$ by a scalar and subtracting $s_t$. What looks like a three-dimensional problem thus costs a two-dimensional one plus a broadcast.

\smalltitle{Feasibility.}
The root of \eqref{\refeq:stateResidual} need not lie inside $\mathcal{S}^{\prime}$: for a given $s_{t-1}$, large $\tau_t$ depresses savings enough that the implied $s_t$ falls below $l_s$, and conversely at small $\tau_t$. We therefore carry an explicit feasibility mask over $\mathcal{S}\times\mathcal{G}_{\tau}$, marking the pairs $(\tau_t, s_{t-1})$ for which a root was found, mask infeasible cells individually, and require at least two feasible $\tau_t$ per state before selecting a maximum there.

\smalltitle{Which derivatives may be taken on the grid.}
Once $s_t(\tau_t,s_{t-1})$ is known, both marginal indirect utilities take the form (level)$^{1-1/\rho}\times$(log-derivative):
\begin{subequations}\label{\refeq:PEEcrraTerms}
\begin{align}
    \frac{d\upsilon_{1,t}^i}{d\tau_t} &= \left(\hat{c}_{1,t}^i\right)^{1-1/\rho}\frac{d\ln\left(\hat{c}_{1,t}^i\right)}{d\tau_t}, &
    \frac{d\upsilon_{2,t}^i}{d\tau_t} &= \left(c_{2,t}^i\right)^{1-1/\rho}\frac{d\ln\left(c_{2,t}^i\right)}{d\tau_t}.
\end{align}
\end{subequations}
For the young households, the discount factor $B_{t+1}^i$ is itself a function of $\tau_t$ through $s_t$, so the factor $(1+B_{t+1}^i)$ in $\upsilon_{1,t}^i$ cannot be held fixed while differentiating. It is convenient to absorb it into the consumption level, defining
\begin{align}\label{\refeq:hatc1i}
    \hat{c}_{1,t}^i \equiv \left(1+B_{t+1}^i\right)^{\frac{1}{1-1/\rho}}\tilde{c}_{1,t}^i \qquad\Longrightarrow\qquad \upsilon_{1,t}^i = \frac{\left(\hat{c}_{1,t}^i\right)^{1-1/\rho}}{1-1/\rho},
\end{align}
so that a single numerical derivative of $\ln(\hat{c}_{1,t}^i)$ picks up both channels at once.\footnote{In implementation $\hat{c}_{1,t}^i$ is never formed as a level. Only $\left(\hat{c}_{1,t}^i\right)^{1-1/\rho} = \left(1+B_{t+1}^i\right)\left(\tilde{c}_{1,t}^i\right)^{1-1/\rho}$ and $\ln\left(\hat{c}_{1,t}^i\right) = \ln\left(1+B_{t+1}^i\right)/(1-1/\rho)+\ln\left(\tilde{c}_{1,t}^i\right)$ are needed, and neither overflows as $\rho\rightarrow 1$, where the literal definition raises $1+B_{t+1}^i$ to a diverging power.}

The two log-derivatives are treated differently. That of $\hat{c}_{1,t}^i$, and as an input to it that of $h_t$, is approximated numerically along the $\tau_t$ dimension of the solution grid, since it depends on $\tau_t$ through the continuation policies and has no closed form. By contrast, $d\ln(c_{2,t}^i)/d\tau_t$ \textbf{must} be taken from its closed form \eqref{\refmodeleq:PEE}: $s_{t-1,i}/s_{t-1}$ varies along the grid, and the policy maker takes it as predetermined, so a grid derivative would include a channel that does not belong in the first order condition. The closed form is the log case's \eqref{\refmodeleq:PEELOG:v2i} with $\partial\ln(\Theta_{h,t})/\partial\tau_t$ replaced by the numerical $d\ln(h_t)/d\tau_t$, which coincide at fixed $s_{t-1}$, so the log and CRRA cases share one formula.

\smalltitle{How the derivatives along $\tau_t$ are taken.}
Each profile along the policy grid, $\ln h_t$ or $\ln\hat{c}_{1,t}^i$ at a given $s_{t-1}$, is fitted by a least-squares cubic spline in $x \equiv \ln(1-\tau_t)$, and the derivative of the fit is mapped back by $dy/d\tau_t = -(dy/dx)/(1-\tau_t)$. Infeasible cells are dropped from the fit and stay undefined. The variable: every profile carries $\ln(1-\tau_t)$ through $\Theta_{h,t}\propto\left((1-\alpha)(1-\tau_t)\right)^{\xi/(1+\alpha\xi)}$, so its derivative in $\tau_t$ grows like $1/(1-\tau_t)$ towards $u$, while in $x$ that part is affine and the cubic reproduces it exactly. The fit: each profile has a kink wherever $s_t(\tau_t)$ crosses a node of the state grid on which the continuation policies are interpolated, and an interpolating spline passes every kink into the root of $z_t$, from where the error compounds backwards through each policy's slope in the state, so that the solution does not settle as $\mathcal{G}_{\tau}$ is refined. The least-squares fit averages over the kinks. Its interior knots are fixed in advance, at the nodes of $\mathcal{G}_{\tau}$ that lie $0.04$ apart in $\tau$, so that the fit is a linear map of the profile and a finer $\mathcal{G}_{\tau}$ puts more nodes under each knot rather than more knots.

\smalltitle{The terminal period, and the log limit.}
The terminal period needs no numerical differentiation: by the chain rule, each term of \eqref{\refmodeleq:terminalPEECRRA} is a consumption level raised to $1-1/\rho$ times the log-derivative already available in closed form from the log case. What is new at $t=T$ is that $B_T^i$ is no longer the primitive $\beta_i$, so the terminal problem $z_T(\tau_T, s_{T-1})$ is state-dependent. The recursion at $t<T$, whose expressions divide by $1-1/\rho$, refuses $\rho = 1$ and directs the caller to \cref{\refsec:PEELOG}. Since the log case is solved by a wholly different route, the two check each other: the CRRA recursion at $\rho$ close to one must reproduce the log path, with a discrepancy that does not shrink as $\rho\rightarrow1$ indicating a defect.

\begin{alg}[Iterative grid search with boundary and feasibility checks]
\label{\refalg:CRRA:grid}
Fix $\bm{\theta}$ and global grids. The following identifies a sequence of policy functions $\tau_t(s_{t-1})$ as well as state policy functions $s_t(\tau_t, s_{t-1})$ that determine the economic state as a function of current taxes and the state along a politico-economic equilibrium path, iterating backwards from the terminal period $t=T$.

\smallskip\noindent\emph{Terminal period.}
Every object entering $z_T$ is closed form in $(\tau_T, s_{T-1})$. On the full grid $\mathcal{S}\times\mathcal{G}_{\tau}$, compute $\Theta_{h,T}$, $h_T$, $B_T^i$, $\Gamma_{s,T-1}$, $s_{T-1,i}/s_{T-1}$, $c_{2,T}^i$, $\tilde{c}_{1,T}^i$ and $d\ln(h_T)/d\tau_T$, assemble $z_T$, and apply the selection rule of \cref{\refsec:RobustRoot} along $\tau$ at each state to obtain $\tau_T(s_{T-1})$ on $\mathcal{S}$. Re-evaluate the auxiliary objects at the selected taxes and construct the interpolants $\tau^T(\cdot)$ and $h^T(\cdot)$ that the previous period will call.

\smallskip\noindent\emph{Iterating backwards for $t<T$.} Given the interpolants $\tau^{t+1}(\cdot)$ and $h^{t+1}(\cdot)$, each period proceeds in four steps.

\begin{enumerate}
    \item \emph{State approximation.} On the Cartesian grid $\mathcal{G}_{\tau}\times\mathcal{S}^{\prime}$ evaluate the forward pass \eqref{\refeq:stateApprox}, retaining $\Theta_{s,t}$ (and $\Theta_{h,t}$, $\Gamma_{s,t}$, $\tau^{t+1}$, $h^{t+1}$, $\bm{B}_{t+1}$ for later use). Form the residual \eqref{\refeq:stateResidual} for every $s_{t-1}\in\mathcal{S}$ by broadcasting, and locate its root along the $s_t$ dimension, giving $s_t(\tau_t, s_{t-1})$ on $\mathcal{G}_{\tau}\times\mathcal{S}$ together with the feasibility mask. This is the only step that touches $\mathcal{S}^{\prime}$, and the only one that calls the continuation interpolants.

    \item \emph{Current-period objects.} With $s_t$ resolved, evaluate on $\mathcal{G}_{\tau}\times\mathcal{S}$: $h_t\left(\Theta_{h,t}, s_{t-1}\right)$, then $R_t$, $\bm{B}_t$, $\Gamma_{s,t-1}\left(\bm{B}_t, \tau_t, \theta_t\right)$ and $s_{t-1,i}/s_{t-1}$, the last two at the previous period's parameter vintage, exactly as in the log case. Then the consumption levels entering \eqref{\refeq:PEEcrraTerms}: $\hat{c}_{1,t}^i$ from \eqref{\refeq:hatc1i} and $c_{2,t}^i$.

    \item \emph{Derivatives and the first order condition.} Approximate $d\ln(h_t)/d\tau_t$ and $d\ln(\hat{c}_{1,t}^i)/d\tau_t$ numerically along $\tau_t$, separately for each $s_{t-1}$, by the least-squares spline in $\ln(1-\tau_t)$ described above; take $d\ln(c_{2,t}^i)/d\tau_t$ from its closed form. Assemble $z_t$ on $\mathcal{G}_{\tau}\times\mathcal{S}$ via \eqref{\refeq:PEEcrraTerms} and the political weights of \eqref{\refeq:LOG:foc}, leaving infeasible cells undefined.

    \item \emph{Selection and reporting.} Apply the selection rule of \cref{\refsec:RobustRoot} along $\tau_t$ for each $s_{t-1}$, skipping infeasible cells, to obtain $\tau_t(s_{t-1})$; the frozen objective of \eqref{\refeq:objectiveProfile} re-uses the numerical derivatives of step 3, so only the retirees' closed-form term is rebuilt per candidate. Re-evaluate steps 1--2 at the selected taxes to obtain the solution on $\mathcal{S}$, and construct the interpolants $\tau_t(s_{t-1})$ and $h_t(s_{t-1})$ that the previous period will call. A light smoothing of the selected $\tau_t(s_{t-1})$ before interpolation removes small kinks inherited from the interpolated continuation policies; its interior knots are fixed in advance (at every $m$-th valid node), so that the smoothed profile is a \emph{linear} map of the profile it smooths and the map from parameters to solution stays continuous for the calibration of \cref{\refsec:calibration}, which differentiates through it.
\end{enumerate}
Steps 1--4 run on the whole grid $\mathcal{S}$ at once, which is why the feasibility mask is carried explicitly.
\end{alg}
